\subsection{MORPHISMS OF LIE GROUPS}

DEFINITION 2. Let G and H be Lie groups. A Lie group morphism of G into H (or simply a morphism of G into H if no confusion can arise) is a mapping of G into H which is a group homomorphism and is analytic. The automorphism group of G is denoted by Aut(G).

The identity mapping of G is a morphism. The composition of two morphisms is a morphism. If \(f: G \to H\) and \(f': H \to G\) are two inverse morphisms, f and \(f'\) are Lie group isomorphisms.

Examples. (1) Let \(G\) be a Lie group. For all \(x \in G\) , \(\operatorname{Int}(x)\) is an automorphism of the Lie group \(G\) .

\begin{enumerate}
\def\labelenumi{(\arabic{enumi})}
\setcounter{enumi}{1}
\item
  Let G be a Lie group. Let \(G^{\nu}\) denote the opposite group to G, with the same manifold structure as G. It is immediate that \(G^{\nu}\) is a Lie group (called the opposite Lie group to G) and that the mapping \(g \mapsto g^{-1}\) is an isomorphism of the Lie group G onto the Lie group \(G^{\nu}\) .
\item
  Let G be a Lie group and E a complete normable space. An analytic linear representation of G on E (or simply a linear representation of G on E when no confusion can arise) is a morphism of the Lie group G into the Lie group \(\mathbf{GL}(\mathrm{E})\) , in other words an analytic mapping \(\pi\) of \(\mathbf{G}\) into \(\mathbf{GL}(\mathrm{E})\) such that \(\pi (gg^{\prime}) = \pi (g)\pi (g^{\prime})\) for \(g\) , \(g^{\prime}\) in \(\mathbf{G}\) . Suppose that E admits a finite basis \((e_1,e_2,\ldots ,e_n)\) over K; let \((e_1^*,e_2^*,\dots ,e_n^*)\) be the dual basis; let \(\rho\) be a homomorphism of the group G into the group \(\mathbf{GL}(\mathrm{E})\) ; then the following conditions are equivalent:
\end{enumerate}

\begin{enumerate}
\def\labelenumi{(\roman{enumi})}
\item
  \(\rho\) is an analytic linear representation;
\item
  for all \(x \in \mathbf{E}\) and \(x' \in \mathbf{E}'\) , the function \(g \mapsto \langle \rho(g)x, x' \rangle\) on \(\mathbf{G}\) is analytic;
\item
  for all \(i\) and \(j\) , the function \(g \mapsto \langle \rho(g)e_i, e_j^*\rangle\) on G is analytic. For the implications (i) \(\Rightarrow\) (ii) \(\Rightarrow\) (iii) are clear. On the other hand, the functions \(u \mapsto \langle ue_i, e_j^*\rangle\) form a coordinate system on \(\mathcal{L}(\mathrm{E})\) ; hence their restrictions to \(\mathbf{GL}(\mathrm{E})\) form a coordinate system on \(\mathbf{GL}(E)\) , whence the implication (iii) \(\Rightarrow\) (i).
\end{enumerate}

Let G be a real Lie group, E a real complete normable space and \(\rho\) a homomorphism of the group G into the group \(\mathbf{GL}(\mathbf{E})\) . We shall see in §8, Theorem 1 that, if \(\rho\) is continuous (when \(\mathbf{GL}(\mathbf{E})\) has the topology derived from the norm on \(\mathrm{L}(\mathbf{E})\) ), then \(\rho\) is analytic. But note that this notion of continuity is different from that considered in Integration, Chapter VIII, §2, Definition 1 (ii) (Exercise 1).

\begin{enumerate}
\def\labelenumi{(\arabic{enumi})}
\setcounter{enumi}{3}
\tightlist
\item
  Let G be a real Lie group and E a complex complete normable space. An analytic linear representation of G on E is a morphism of G into the underlying real Lie group of \(\mathbf{GL}(\mathrm{E})\) .
\end{enumerate}

PROPOSITION 4. Let G and H be Lie groups and f a homomorphism of the group G into the group H. For f to be analytic, it is necessary and sufficient that there exist a non-empty open subset U of G such that \(f \mid U\) is analytic.

The condition is obviously necessary. Suppose that it holds. For all \(x_0 \in G\) , \(f(x_0x) = f(x_0)f(x)\) for all \(x \in U\) and hence \(f|_{x_0U}\) is analytic. But the sets \(x_0U\) , where \(x_0 \in G\) , form an open covering of \(G\) .

Remark. If \(f\) is an immersion at \(e\) (resp. a submersion at \(e\) ), clearly \(f\) is an immersion (resp. a submersion).

